%%%%%%%%%%%%%%%%%%%%%%%%%%%%%%%%%%%%%%%%%%%%%%%%%%%%%%%%%%%%%%%%%%%%%%%%%%%%%%%%
% CHAPTER 10
%%%%%%%%%%%%%%%%%%%%%%%%%%%%%%%%%%%%%%%%%%%%%%%%%%%%%%%%%%%%%%%%%%%%%%%%%%%%%%%%

\chapter{Fast Fourier Transform}

\begin{tcolorbox}[colback=blue!5!white,colframe=blue!70!black,title=Learning Objectives]

After studying this chapter, the reader will be able to

\begin{itemize}
\item Understand the need for FFT.
\item Compare DFT and FFT computational complexity.
\item Derive the radix-2 decimation-in-time FFT.
\item Understand butterfly operations.
\item Perform a 4-point FFT manually.
\item Understand bit-reversal ordering.
\end{itemize}

\end{tcolorbox}

%%%%%%%%%%%%%%%%%%%%%%%%%%%%%%%%%%%%%%%%%%%%%%%%%%%%%%%%%%%%%%%%%%%%%%%%%%%%%%%%

\section{Introduction}

The DFT of an \(N\)-point sequence is

\[
X[k]=
\sum_{n=0}^{N-1}
x[n]W_N^{kn},
\qquad
W_N=e^{-j2\pi/N}.
\]

A direct computation requires

\[
N^2
\]

complex multiplications.

For large \(N\), this becomes computationally expensive.

%%%%%%%%%%%%%%%%%%%%%%%%%%%%%%%%%%%%%%%%%%%%%%%%%%%%%%%%%%%%%%%%%%%%%%%%%%%%%%%%

\section{Need for FFT}

The \textbf{Fast Fourier Transform (FFT)} is an efficient algorithm for computing
the DFT.

%%%%%%%%%%%%%%%%%%%%%%%%%%%%%%%%%%%%%%%%%%%%%%%%%%%%%%%%%%%%%%%%%%%%%%%%%%%%%%%%

\subsection{Complexity Comparison}

\begin{table}[H]
\centering
\caption{DFT vs FFT}
\begin{tabular}{|c|c|c|}
\hline
\textbf{\(N\)} & \textbf{DFT multiplications} & \textbf{FFT multiplications}\\
\hline
8 & 64 & 12\\
16 & 256 & 32\\
64 & 4096 & 192\\
1024 & 1,048,576 & 5120\\
\hline
\end{tabular}
\end{table}

%%%%%%%%%%%%%%%%%%%%%%%%%%%%%%%%%%%%%%%%%%%%%%%%%%%%%%%%%%%%%%%%%%%%%%%%%%%%%%%%

\subsection{FFT Complexity}

For radix-2 FFT,

\[
\boxed{
\text{Complexity}=N\log_2N
}
\]

This reduction makes real-time DSP practical.

%%%%%%%%%%%%%%%%%%%%%%%%%%%%%%%%%%%%%%%%%%%%%%%%%%%%%%%%%%%%%%%%%%%%%%%%%%%%%%%%

\section{Divide-and-Conquer Idea}

The FFT divides the \(N\)-point DFT into smaller DFTs.

%%%%%%%%%%%%%%%%%%%%%%%%%%%%%%%%%%%%%%%%%%%%%%%%%%%%%%%%%%%%%%%%%%%%%%%%%%%%%%%%

\subsection{Even-Odd Decomposition}

Separate the sequence into

\begin{itemize}
\item even-indexed samples,
\item odd-indexed samples.
\end{itemize}

\[
x_e[r]=x[2r]
\]

\[
x_o[r]=x[2r+1]
\]

where

\[
r=0,1,\ldots,\frac{N}{2}-1.
\]

%%%%%%%%%%%%%%%%%%%%%%%%%%%%%%%%%%%%%%%%%%%%%%%%%%%%%%%%%%%%%%%%%%%%%%%%%%%%%%%%

\section{Derivation of Radix-2 DIT FFT}

Start with

\[
X[k]=
\sum_{n=0}^{N-1}x[n]W_N^{kn}.
\]

Split the sum into even and odd terms:

\[
X[k]=
\sum_{r=0}^{N/2-1}x[2r]W_N^{k(2r)}
+
\sum_{r=0}^{N/2-1}x[2r+1]W_N^{k(2r+1)}.
\]

Using

\[
W_N^{2k}=W_{N/2}^k,
\]

we get

\[
X[k]=
\sum_{r=0}^{N/2-1}x_e[r]W_{N/2}^{kr}
+
W_N^k
\sum_{r=0}^{N/2-1}x_o[r]W_{N/2}^{kr}.
\]

Define

\[
E[k]=
\sum_{r=0}^{N/2-1}x_e[r]W_{N/2}^{kr},
\]

\[
O[k]=
\sum_{r=0}^{N/2-1}x_o[r]W_{N/2}^{kr}.
\]

Then

\[
\boxed{
X[k]=E[k]+W_N^kO[k]
}
\]

and

\[
\boxed{
X[k+N/2]=E[k]-W_N^kO[k]
}
\]

for

\[
k=0,1,\ldots,\frac{N}{2}-1.
\]

%%%%%%%%%%%%%%%%%%%%%%%%%%%%%%%%%%%%%%%%%%%%%%%%%%%%%%%%%%%%%%%%%%%%%%%%%%%%%%%%

\section{Butterfly Operation}

The pair of equations above is implemented by a \textbf{butterfly}.

\begin{center}
\begin{tikzpicture}[>=Latex]
\node at (0,1) {$E[k]$};
\node at (0,-1) {$O[k]$};

\draw[->] (0.6,1) -- (2.2,0.4);
\draw[->] (0.6,-1) -- (2.2,-0.4);

\draw[->] (0.6,1) -- (2.2,-0.4);
\draw[->] (0.6,-1) -- (2.2,0.4);

\node[right] at (2.4,0.4) {$X[k]$};
\node[right] at (2.4,-0.4) {$X[k+N/2]$};

\node at (1.4,-1.4) {$W_N^k$};
\end{tikzpicture}
\end{center}

The butterfly performs

\[
\boxed{
\begin{aligned}
X[k] &= E[k]+W_N^kO[k]\\
X[k+N/2] &= E[k]-W_N^kO[k]
\end{aligned}
}
\]

%%%%%%%%%%%%%%%%%%%%%%%%%%%%%%%%%%%%%%%%%%%%%%%%%%%%%%%%%%%%%%%%%%%%%%%%%%%%%%%%

\section{4-Point FFT}

Let

\[
x[n]=\{1,2,3,4\}.
\]

%%%%%%%%%%%%%%%%%%%%%%%%%%%%%%%%%%%%%%%%%%%%%%%%%%%%%%%%%%%%%%%%%%%%%%%%%%%%%%%%

\subsection{Step 1: Separate Even and Odd Samples}

Even sequence:

\[
x_e=\{1,3\}
\]

Odd sequence:

\[
x_o=\{2,4\}.
\]

%%%%%%%%%%%%%%%%%%%%%%%%%%%%%%%%%%%%%%%%%%%%%%%%%%%%%%%%%%%%%%%%%%%%%%%%%%%%%%%%

\section{2-Point DFT of Even Sequence}

For a 2-point DFT,

\[
E[0]=1+3=4
\]

\[
E[1]=1-3=-2.
\]

%%%%%%%%%%%%%%%%%%%%%%%%%%%%%%%%%%%%%%%%%%%%%%%%%%%%%%%%%%%%%%%%%%%%%%%%%%%%%%%%

\section{2-Point DFT of Odd Sequence}

\[
O[0]=2+4=6
\]

\[
O[1]=2-4=-2.
\]

%%%%%%%%%%%%%%%%%%%%%%%%%%%%%%%%%%%%%%%%%%%%%%%%%%%%%%%%%%%%%%%%%%%%%%%%%%%%%%%%

\section{Twiddle Factors}

For \(N=4\),

\[
W_4=e^{-j\pi/2}=-j.
\]

\[
W_4^0=1,
\qquad
W_4^1=-j.
\]

%%%%%%%%%%%%%%%%%%%%%%%%%%%%%%%%%%%%%%%%%%%%%%%%%%%%%%%%%%%%%%%%%%%%%%%%%%%%%%%%

\section{Combine for \(k=0\)}

\[
X[0]=E[0]+W_4^0O[0]
\]

\[
=4+6=10.
\]

\[
X[2]=E[0]-W_4^0O[0]
\]

\[
=4-6=-2.
\]

%%%%%%%%%%%%%%%%%%%%%%%%%%%%%%%%%%%%%%%%%%%%%%%%%%%%%%%%%%%%%%%%%%%%%%%%%%%%%%%%

\section{Combine for \(k=1\)}

\[
X[1]=E[1]+W_4^1O[1]
\]

\[
=-2+(-j)(-2)
=-2+2j.
\]

\[
X[3]=E[1]-W_4^1O[1]
\]

\[
=-2-2j.
\]

%%%%%%%%%%%%%%%%%%%%%%%%%%%%%%%%%%%%%%%%%%%%%%%%%%%%%%%%%%%%%%%%%%%%%%%%%%%%%%%%

\section{Final Result}

\[
\boxed{
X[k]=
\{10,\,-2+2j,\,-2,\,-2-2j\}
}
\]

%%%%%%%%%%%%%%%%%%%%%%%%%%%%%%%%%%%%%%%%%%%%%%%%%%%%%%%%%%%%%%%%%%%%%%%%%%%%%%%%

\section{Verification}

The direct DFT gives the same result, confirming the correctness of the FFT
computation.

%%%%%%%%%%%%%%%%%%%%%%%%%%%%%%%%%%%%%%%%%%%%%%%%%%%%%%%%%%%%%%%%%%%%%%%%%%%%%%%%

\section{FFT Stages}

For radix-2 FFT,

\[
\boxed{
\text{Number of stages}=\log_2N
}
\]

Examples:

\begin{table}[H]
\centering
\caption{FFT stages}
\begin{tabular}{|c|c|}
\hline
\textbf{\(N\)} & \textbf{Stages}\\
\hline
4 & 2\\
8 & 3\\
16 & 4\\
32 & 5\\
64 & 6\\
\hline
\end{tabular}
\end{table}

%%%%%%%%%%%%%%%%%%%%%%%%%%%%%%%%%%%%%%%%%%%%%%%%%%%%%%%%%%%%%%%%%%%%%%%%%%%%%%%%

\section{Butterflies per Stage}

Each stage contains

\[
\boxed{
\frac{N}{2}
}
\]

butterflies.

Total butterflies:

\[
\boxed{
\frac{N}{2}\log_2N
}
\]

%%%%%%%%%%%%%%%%%%%%%%%%%%%%%%%%%%%%%%%%%%%%%%%%%%%%%%%%%%%%%%%%%%%%%%%%%%%%%%%%

\section{Bit-Reversal Ordering}

In radix-2 DIT FFT, the input must be arranged in bit-reversed order.

%%%%%%%%%%%%%%%%%%%%%%%%%%%%%%%%%%%%%%%%%%%%%%%%%%%%%%%%%%%%%%%%%%%%%%%%%%%%%%%%

\subsection{Example for \(N=8\)}

\begin{table}[H]
\centering
\caption{Bit-reversal table}
\begin{tabular}{|c|c|c|}
\hline
\textbf{Decimal} & \textbf{Binary} & \textbf{Reversed}\\
\hline
0 & 000 & 000 (0)\\
1 & 001 & 100 (4)\\
2 & 010 & 010 (2)\\
3 & 011 & 110 (6)\\
4 & 100 & 001 (1)\\
5 & 101 & 101 (5)\\
6 & 110 & 011 (3)\\
7 & 111 & 111 (7)\\
\hline
\end{tabular}
\end{table}

The input order becomes

\[
\boxed{
\{0,4,2,6,1,5,3,7\}
}
\]

%%%%%%%%%%%%%%%%%%%%%%%%%%%%%%%%%%%%%%%%%%%%%%%%%%%%%%%%%%%%%%%%%%%%%%%%%%%%%%%%

\section{Why Bit Reversal?}

Bit reversal allows the FFT to be computed \textbf{in place}, meaning the input
array is overwritten by intermediate results and finally by the output.

This reduces memory usage significantly.

%%%%%%%%%%%%%%%%%%%%%%%%%%%%%%%%%%%%%%%%%%%%%%%%%%%%%%%%%%%%%%%%%%%%%%%%%%%%%%%%

\section{Flow Graph of 4-Point FFT}

\begin{center}
\begin{tikzpicture}[x=1.6cm,y=1cm,>=Latex]
\node (x0) at (0,3) {$x[0]$};
\node (x2) at (0,2) {$x[2]$};
\node (x1) at (0,1) {$x[1]$};
\node (x3) at (0,0) {$x[3]$};

\draw (x0) -- (1,2.5);
\draw (x2) -- (1,2.5);
\draw (x0) -- (1,1.5);
\draw (x2) -- (1,1.5);

\draw (x1) -- (1,0.5);
\draw (x3) -- (1,0.5);
\draw (x1) -- (1,-0.5);
\draw (x3) -- (1,-0.5);

\draw (1,2.5) -- (2,2);
\draw (1,1.5) -- (2,1);
\draw (1,0.5) -- (2,1);
\draw (1,-0.5) -- (2,0);

\node[right] at (2,2) {$X[0]$};
\node[right] at (2,1) {$X[1]$};
\node[right] at (2,0) {$X[2],X[3]$};
\end{tikzpicture}
\end{center}

%%%%%%%%%%%%%%%%%%%%%%%%%%%%%%%%%%%%%%%%%%%%%%%%%%%%%%%%%%%%%%%%%%%%%%%%%%%%%%%%

\section{Computational Saving}

For \(N=1024\),

Direct DFT:

\[
1024^2=1,048,576
\]

multiplications.

FFT:

\[
1024\log_2(1024)=1024\times10=10,240
\]

operations.

This is approximately

\[
\boxed{
100\times \text{ fewer operations}
}
\]

%%%%%%%%%%%%%%%%%%%%%%%%%%%%%%%%%%%%%%%%%%%%%%%%%%%%%%%%%%%%%%%%%%%%%%%%%%%%%%%%

\section{Quick Check}

\begin{enumerate}
\item \(W_4=\boxed{-j}\).
\item 8-point FFT has \(\boxed{3}\) stages.
\item Butterflies per stage for \(N=16\): \(\boxed{8}\).
\item Total butterflies for \(N=16\): \(\boxed{32}\).
\item Bit-reversal of binary \(101\) is \(\boxed{101}\) (decimal 5).
\end{enumerate}

%%%%%%%%%%%%%%%%%%%%%%%%%%%%%%%%%%%%%%%%%%%%%%%%%%%%%%%%%%%%%%%%%%%%%%%%%%%%%%%%

\begin{tcolorbox}[title=One-Minute Revision,colback=blue!5]

\[
X[k]=
\sum_{n=0}^{N-1}x[n]W_N^{kn}
\]

\[
W_N=e^{-j2\pi/N}
\]

Radix-2 DIT:

\[
X[k]=E[k]+W_N^kO[k]
\]

\[
X[k+N/2]=E[k]-W_N^kO[k]
\]

Stages:

\[
\log_2N
\]

Butterflies per stage:

\[
\frac{N}{2}
\]

Total butterflies:

\[
\frac{N}{2}\log_2N
\]

Bit-reversal is required for DIT FFT input ordering.

FFT complexity:

\[
N\log_2N
\]

\end{tcolorbox}
%%%%%%%%%%%%%%%%%%%%%%%%%%%%%%%%%%%%%%%%%%%%%%%%%%%%%%%%%%%%%%%%%%%%%%%%%%%%%%%%
%%%%%%%%%%%%%%%%%%%%%%%%%%%%%%%%%%%%%%%%%%%%%%%%%%%%%%%%%%%%%%%%%%%%%%%%%%%%%%%%
\section{Radix-2 Decimation-in-Frequency (DIF) FFT}

In the DIT algorithm, the input is divided into even and odd samples.

In the \textbf{DIF algorithm}, the frequency-domain outputs are divided into
even and odd frequency indices.

Start with

\[
X[k]=
\sum_{n=0}^{N-1}x[n]W_N^{kn}.
\]

Split the sequence into two halves:

\[
X[k]=
\sum_{n=0}^{N/2-1}x[n]W_N^{kn}
+
\sum_{n=0}^{N/2-1}x[n+N/2]W_N^{k(n+N/2)}.
\]

Using

\[
W_N^{kN/2}=(-1)^k,
\]

we obtain

\[
X[k]=
\sum_{n=0}^{N/2-1}
\left[
x[n]+(-1)^k x[n+N/2]
\right]
W_{N/2}^{kn}.
\]

%%%%%%%%%%%%%%%%%%%%%%%%%%%%%%%%%%%%%%%%%%%%%%%%%%%%%%%%%%%%%%%%%%%%%%%%%%%%%%%%

\section{DIF Butterfly Equations}

For

\[
n=0,1,\ldots,\frac{N}{2}-1,
\]

compute

\[
\boxed{
u[n]=x[n]+x[n+N/2]
}
\]

\[
\boxed{
v[n]=
\left(x[n]-x[n+N/2]\right)W_N^n
}
\]

Then perform two \((N/2)\)-point FFTs on \(u[n]\) and \(v[n]\).

%%%%%%%%%%%%%%%%%%%%%%%%%%%%%%%%%%%%%%%%%%%%%%%%%%%%%%%%%%%%%%%%%%%%%%%%%%%%%%%%

\section{DIT vs DIF}

\begin{table}[H]
\centering
\caption{DIT and DIF comparison}
\begin{tabular}{|p{4cm}|p{4cm}|p{4cm}|}
\hline
\textbf{Feature} & \textbf{DIT} & \textbf{DIF}\\
\hline
Decimation & Time & Frequency\\
\hline
Input order & Bit-reversed & Normal\\
\hline
Output order & Normal & Bit-reversed\\
\hline
Twiddle multiplication & Before butterfly output & After subtraction\\
\hline
Most common in textbooks & Yes & Yes\\
\hline
\end{tabular}
\end{table}

%%%%%%%%%%%%%%%%%%%%%%%%%%%%%%%%%%%%%%%%%%%%%%%%%%%%%%%%%%%%%%%%%%%%%%%%%%%%%%%%

\section{In-Place Computation}

FFT algorithms can overwrite the input array with intermediate results.

For example,

\[
x[0],x[1],x[2],x[3]
\]

are updated stage by stage until they contain

\[
X[0],X[1],X[2],X[3].
\]

This is called \textbf{in-place computation} and saves memory.

%%%%%%%%%%%%%%%%%%%%%%%%%%%%%%%%%%%%%%%%%%%%%%%%%%%%%%%%%%%%%%%%%%%%%%%%%%%%%%%%

\section{Inverse FFT (IFFT)}

The IDFT is

\[
x[n]=
\frac{1}{N}
\sum_{k=0}^{N-1}X[k]e^{j2\pi kn/N}.
\]

Notice that the sign in the exponent is opposite to the DFT.

%%%%%%%%%%%%%%%%%%%%%%%%%%%%%%%%%%%%%%%%%%%%%%%%%%%%%%%%%%%%%%%%%%%%%%%%%%%%%%%%

\subsection{IFFT Using FFT}

The IFFT can be computed using a forward FFT:

\[
\boxed{
x[n]=
\frac{1}{N}
\left(
\text{FFT}\{X^*[k]\}
\right)^*
}
\]

Procedure:

\begin{enumerate}
\item Take the complex conjugate of \(X[k]\).
\item Compute the FFT.
\item Take the conjugate of the result.
\item Divide by \(N\).
\end{enumerate}

%%%%%%%%%%%%%%%%%%%%%%%%%%%%%%%%%%%%%%%%%%%%%%%%%%%%%%%%%%%%%%%%%%%%%%%%%%%%%%%%

\section{Example: IFFT}

Let

\[
X[k]=\{4,0,0,0\}.
\]

Applying the IFFT,

\[
x[n]=\{1,1,1,1\}.
\]

%%%%%%%%%%%%%%%%%%%%%%%%%%%%%%%%%%%%%%%%%%%%%%%%%%%%%%%%%%%%%%%%%%%%%%%%%%%%%%%%

\section{FFT-Based Linear Convolution}

Direct convolution of long sequences is computationally expensive.

%%%%%%%%%%%%%%%%%%%%%%%%%%%%%%%%%%%%%%%%%%%%%%%%%%%%%%%%%%%%%%%%%%%%%%%%%%%%%%%%

\subsection{Procedure}

For sequences of lengths \(L_x\) and \(L_h\):

\begin{enumerate}
\item Choose
\[
N\ge L_x+L_h-1.
\]

\item Zero-pad both sequences to length \(N\).

\item Compute FFTs:
\[
X[k]=\text{FFT}\{x[n]\},
\qquad
H[k]=\text{FFT}\{h[n]\}.
\]

\item Multiply:
\[
Y[k]=X[k]H[k].
\]

\item Compute
\[
y[n]=\text{IFFT}\{Y[k]\}.
\]
\end{enumerate}

%%%%%%%%%%%%%%%%%%%%%%%%%%%%%%%%%%%%%%%%%%%%%%%%%%%%%%%%%%%%%%%%%%%%%%%%%%%%%%%%

\section{Example: FFT Convolution}

Let

\[
x[n]=\{1,2\},
\qquad
h[n]=\{1,1\}.
\]

Output length:

\[
L=2+2-1=3.
\]

Choose \(N=4\).

%%%%%%%%%%%%%%%%%%%%%%%%%%%%%%%%%%%%%%%%%%%%%%%%%%%%%%%%%%%%%%%%%%%%%%%%%%%%%%%%

\subsection{Zero-Padded Sequences}

\[
x[n]=\{1,2,0,0\},
\]

\[
h[n]=\{1,1,0,0\}.
\]

%%%%%%%%%%%%%%%%%%%%%%%%%%%%%%%%%%%%%%%%%%%%%%%%%%%%%%%%%%%%%%%%%%%%%%%%%%%%%%%%

\subsection{Result}

The IFFT gives

\[
\boxed{
y[n]=\{1,3,2,0\}.
}
\]

Ignoring the trailing zero,

\[
\boxed{
y[n]=\{1,3,2\}
}
\]

which is the correct linear convolution.

%%%%%%%%%%%%%%%%%%%%%%%%%%%%%%%%%%%%%%%%%%%%%%%%%%%%%%%%%%%%%%%%%%%%%%%%%%%%%%%%

\section{Block Convolution}

For very long signals, convolution is performed block by block.

Two standard methods are

\begin{itemize}
\item Overlap-Add (OLA),
\item Overlap-Save (OLS).
\end{itemize}

%%%%%%%%%%%%%%%%%%%%%%%%%%%%%%%%%%%%%%%%%%%%%%%%%%%%%%%%%%%%%%%%%%%%%%%%%%%%%%%%

\section{Overlap-Add Method}

%%%%%%%%%%%%%%%%%%%%%%%%%%%%%%%%%%%%%%%%%%%%%%%%%%%%%%%%%%%%%%%%%%%%%%%%%%%%%%%%

\subsection{Idea}

\begin{enumerate}
\item Divide the input into non-overlapping blocks.
\item Convolve each block with the filter using FFT.
\item Add the overlapping output portions.
\end{enumerate}

%%%%%%%%%%%%%%%%%%%%%%%%%%%%%%%%%%%%%%%%%%%%%%%%%%%%%%%%%%%%%%%%%%%%%%%%%%%%%%%%

\subsection{Block Diagram}

\begin{center}
\begin{tikzpicture}[node distance=2.2cm,>=Latex]
\node[draw,minimum width=2.4cm,minimum height=0.9cm] (blk) {Input blocks};
\node[draw,minimum width=2.6cm,minimum height=0.9cm,right of=blk] (fft) {FFT convolution};
\node[draw,minimum width=2.4cm,minimum height=0.9cm,right of=fft] (add) {Overlap \& Add};

\draw[->] (blk) -- (fft);
\draw[->] (fft) -- (add);
\end{tikzpicture}
\end{center}

%%%%%%%%%%%%%%%%%%%%%%%%%%%%%%%%%%%%%%%%%%%%%%%%%%%%%%%%%%%%%%%%%%%%%%%%%%%%%%%%

\subsection{Advantage}

Simple to understand and implement.

%%%%%%%%%%%%%%%%%%%%%%%%%%%%%%%%%%%%%%%%%%%%%%%%%%%%%%%%%%%%%%%%%%%%%%%%%%%%%%%%

\subsection{Disadvantage}

Requires addition of overlapping segments.

%%%%%%%%%%%%%%%%%%%%%%%%%%%%%%%%%%%%%%%%%%%%%%%%%%%%%%%%%%%%%%%%%%%%%%%%%%%%%%%%

\section{Overlap-Save Method}

%%%%%%%%%%%%%%%%%%%%%%%%%%%%%%%%%%%%%%%%%%%%%%%%%%%%%%%%%%%%%%%%%%%%%%%%%%%%%%%%

\subsection{Idea}

\begin{enumerate}
\item Consecutive blocks overlap by \(L_h-1\) samples.
\item Perform circular convolution using FFT.
\item Discard the corrupted samples at the beginning.
\item Keep the remaining valid samples.
\end{enumerate}

%%%%%%%%%%%%%%%%%%%%%%%%%%%%%%%%%%%%%%%%%%%%%%%%%%%%%%%%%%%%%%%%%%%%%%%%%%%%%%%%

\subsection{Advantage}

No overlap addition is required.

%%%%%%%%%%%%%%%%%%%%%%%%%%%%%%%%%%%%%%%%%%%%%%%%%%%%%%%%%%%%%%%%%%%%%%%%%%%%%%%%

\subsection{Disadvantage}

Input block management is slightly more complex.

%%%%%%%%%%%%%%%%%%%%%%%%%%%%%%%%%%%%%%%%%%%%%%%%%%%%%%%%%%%%%%%%%%%%%%%%%%%%%%%%

\section{OLA vs OLS}

\begin{table}[H]
\centering
\caption{Comparison of block convolution methods}
\begin{tabular}{|p{4cm}|p{4cm}|p{4cm}|}
\hline
\textbf{Feature} & \textbf{OLA} & \textbf{OLS}\\
\hline
Input blocks & Non-overlapping & Overlapping\\
\hline
Output processing & Add overlaps & Discard corrupted samples\\
\hline
Extra additions & Yes & No\\
\hline
Memory requirement & Moderate & Moderate\\
\hline
Practical use & Common & Very common\\
\hline
\end{tabular}
\end{table}

%%%%%%%%%%%%%%%%%%%%%%%%%%%%%%%%%%%%%%%%%%%%%%%%%%%%%%%%%%%%%%%%%%%%%%%%%%%%%%%%

\section{Computational Complexity}

For radix-2 FFT:

%%%%%%%%%%%%%%%%%%%%%%%%%%%%%%%%%%%%%%%%%%%%%%%%%%%%%%%%%%%%%%%%%%%%%%%%%%%%%%%%

\subsection{Complex Multiplications}

\[
\boxed{
\frac{N}{2}\log_2N
}
\]

%%%%%%%%%%%%%%%%%%%%%%%%%%%%%%%%%%%%%%%%%%%%%%%%%%%%%%%%%%%%%%%%%%%%%%%%%%%%%%%%

\subsection{Complex Additions}

\[
\boxed{
N\log_2N
}
\]

%%%%%%%%%%%%%%%%%%%%%%%%%%%%%%%%%%%%%%%%%%%%%%%%%%%%%%%%%%%%%%%%%%%%%%%%%%%%%%%%

\section{Example: 8-Point FFT}

\[
N=8,
\qquad
\log_2 8=3.
\]

Multiplications:

\[
\frac{8}{2}\times3=12.
\]

Additions:

\[
8\times3=24.
\]

%%%%%%%%%%%%%%%%%%%%%%%%%%%%%%%%%%%%%%%%%%%%%%%%%%%%%%%%%%%%%%%%%%%%%%%%%%%%%%%%

\section{Example: 16-Point FFT}

\[
N=16,
\qquad
\log_2 16=4.
\]

Multiplications:

\[
\frac{16}{2}\times4=32.
\]

Additions:

\[
16\times4=64.
\]

%%%%%%%%%%%%%%%%%%%%%%%%%%%%%%%%%%%%%%%%%%%%%%%%%%%%%%%%%%%%%%%%%%%%%%%%%%%%%%%%

\section{Solved Example 10.1}

How many stages are required for a 32-point radix-2 FFT?

%%%%%%%%%%%%%%%%%%%%%%%%%%%%%%%%%%%%%%%%%%%%%%%%%%%%%%%%%%%%%%%%%%%%%%%%%%%%%%%%

\subsection{Solution}

\[
\log_2 32=5.
\]

\[
\boxed{
5\text{ stages}
}
\]

%%%%%%%%%%%%%%%%%%%%%%%%%%%%%%%%%%%%%%%%%%%%%%%%%%%%%%%%%%%%%%%%%%%%%%%%%%%%%%%%

\section{Solved Example 10.2}

How many butterflies are required for a 32-point FFT?

%%%%%%%%%%%%%%%%%%%%%%%%%%%%%%%%%%%%%%%%%%%%%%%%%%%%%%%%%%%%%%%%%%%%%%%%%%%%%%%%

\subsection{Solution}

Butterflies per stage:

\[
\frac{32}{2}=16.
\]

Total butterflies:

\[
16\times5=80.
\]

\[
\boxed{
80\text{ butterflies}
}
\]

%%%%%%%%%%%%%%%%%%%%%%%%%%%%%%%%%%%%%%%%%%%%%%%%%%%%%%%%%%%%%%%%%%%%%%%%%%%%%%%%

\section{Solved Example 10.3}

Find the bit-reversed order for \(N=8\).

Using the bit-reversal table:

\[
0\rightarrow0,
\quad
1\rightarrow4,
\quad
2\rightarrow2,
\quad
3\rightarrow6,
\]

\[
4\rightarrow1,
\quad
5\rightarrow5,
\quad
6\rightarrow3,
\quad
7\rightarrow7.
\]

Hence,

\[
\boxed{
\{0,4,2,6,1,5,3,7\}
}
\]

%%%%%%%%%%%%%%%%%%%%%%%%%%%%%%%%%%%%%%%%%%%%%%%%%%%%%%%%%%%%%%%%%%%%%%%%%%%%%%%%

\section{Solved Example 10.4}

A 64-point FFT is used. Find the number of complex multiplications.

%%%%%%%%%%%%%%%%%%%%%%%%%%%%%%%%%%%%%%%%%%%%%%%%%%%%%%%%%%%%%%%%%%%%%%%%%%%%%%%%

\subsection{Solution}

\[
\frac{64}{2}\log_2 64
=
32\times6
=
192.
\]

\[
\boxed{
192
}
\]

%%%%%%%%%%%%%%%%%%%%%%%%%%%%%%%%%%%%%%%%%%%%%%%%%%%%%%%%%%%%%%%%%%%%%%%%%%%%%%%%

\section{Solved Example 10.5}

A sequence of length 100 is convolved with a filter of length 25 using FFT.

Minimum FFT size:

\[
N\ge100+25-1=124.
\]

Choose the next power of two:

\[
\boxed{
N=128
}
\]

%%%%%%%%%%%%%%%%%%%%%%%%%%%%%%%%%%%%%%%%%%%%%%%%%%%%%%%%%%%%%%%%%%%%%%%%%%%%%%%%

\section{Important GATE Points}

\begin{enumerate}
\item FFT is an algorithm, not a new transform.
\item Radix-2 FFT requires \(N=2^m\).
\item DIT: bit-reversed input, normal output.
\item DIF: normal input, bit-reversed output.
\item IFFT can be computed using FFT and conjugation.
\item FFT-based convolution requires zero padding.
\item OLA adds overlaps; OLS discards corrupted samples.
\item Complexity is \(N\log_2N\), not \(N^2\).
\end{enumerate}

%%%%%%%%%%%%%%%%%%%%%%%%%%%%%%%%%%%%%%%%%%%%%%%%%%%%%%%%%%%%%%%%%%%%%%%%%%%%%%%%

\section{Chapter Summary}

In this chapter, we derived the radix-2 DIT FFT and introduced the DIF FFT,
butterfly operations, bit-reversal ordering and in-place computation.

We also studied the IFFT, FFT-based linear convolution and the practical block
convolution methods OLA and OLS.

These techniques are used in spectrum analyzers, communication systems, audio
processing, radar, image processing and real-time embedded DSP systems.

%%%%%%%%%%%%%%%%%%%%%%%%%%%%%%%%%%%%%%%%%%%%%%%%%%%%%%%%%%%%%%%%%%%%%%%%%%%%%%%%

\begin{tcolorbox}[title=One-Minute Revision,colback=blue!5]

FFT complexity:

\[
N\log_2N
\]

DIT butterfly:

\[
X[k]=E[k]+W_N^kO[k]
\]

\[
X[k+N/2]=E[k]-W_N^kO[k]
\]

DIF butterfly:

\[
u[n]=x[n]+x[n+N/2]
\]

\[
v[n]=(x[n]-x[n+N/2])W_N^n
\]

Stages:

\[
\log_2N
\]

Butterflies:

\[
\frac{N}{2}\log_2N
\]

IFFT:

\[
x[n]=
\frac{1}{N}
\left(
\text{FFT}\{X^*[k]\}
\right)^*
\]

Linear convolution:

\[
N\ge L_x+L_h-1
\]

OLA: add overlaps.

OLS: discard corrupted samples.

\end{tcolorbox}
%%%%%%%%%%%%%%%%%%%%%%%%%%%%%%%%%%%%%%%%%%%%%%%%%%%%%%%%%%%%%%%%%%%%%%%%%%%%%%%%